% Heartwood: Publicly Verifiable Proof of Semantic Effort for Black-Box LLM Endpoints
%
% Build:  pdflatex heartwood.tex  (twice, for references)
%
\documentclass[11pt]{article}

\usepackage[margin=1in]{geometry}
\usepackage[T1]{fontenc}
\usepackage{lmodern}
\usepackage{amsmath,amssymb,amsthm}
\usepackage{booktabs}
\usepackage{url}
\usepackage[hidelinks]{hyperref}
\usepackage{microtype}

\newtheorem{theorem}{Theorem}
\newtheorem{proposition}{Proposition}
\theoremstyle{definition}
\newtheorem{definition}{Definition}

\title{Heartwood: Publicly Verifiable Proof of Semantic Effort\\
for Black-Box LLM Endpoints}

\author{George Tejada\thanks{Independent researcher. Reference implementation,
receipts, and all evaluation artefacts are MIT-licensed and public at
\url{https://github.com/zcashsensei/heartwood}.}}

\date{August 2026}

\begin{document}
\maketitle

\begin{abstract}
A customer calling a hosted language model cannot tell how much computation the
provider actually spent on their request. Existing verification layers each bind
a different object and none binds effort: trusted-execution attestation binds the
hardware and the loaded image, zero-knowledge proofs of inference bind the
architecture and committed weights, and signed-response provenance binds a
signature chain. The Hollow-LLM attack establishes that this is not an
engineering gap but a structural one --- ``proof of correct inference is not
proof of large-model execution'' --- because a prover can satisfy the circuit
while collapsing the effective computation. Behavioural auditing does measure
capability, but the state of the art states its own limitation plainly: such
audits are auditor-centric, with no mechanism for independent verification and
no use of cryptographic commitment or public randomness.

We close the gap between these two literatures. Our observation is that
\emph{a correct answer to a problem requiring $N$ sequential steps is itself
evidence that $N$ steps were spent}: effort can be bound semantically, in the
output, where cryptography provably cannot bind it in the computation. We make
that evidence transferable by committing a challenge pool before a public
randomness beacon selects which items are used, and by accumulating evidence as
an anytime-valid test supermartingale, so neither cherry-picking nor optional
stopping can bias the result. The output is a \emph{receipt} that any third
party recomputes offline, with no contact with the auditor or the provider and
no cooperation from the provider at any point.

On production frontier APIs, Heartwood detects a full effort skim on Claude
Opus~5 in $4$ queries and on Claude Haiku~4.5 in $4$ queries at
$\alpha = 0.01$, with no false positive across $105$ honest queries. We report a
property we did not anticipate: the audit becomes \emph{cheaper} on more capable
models, from $31$ queries on a 2B local model to $4$ on frontier models, because
a capable model fails harder when its compute is removed. We also report two
negative results that constrain the method --- disabling a reasoning parameter
is not by itself an effort skim, and a degradation milder than the declared
tolerance is not detected --- and we publish the receipts for both.
\end{abstract}

\section{Introduction}

Hosted language models are sold as a named capability. A customer pays for a
particular model at a particular reasoning tier and receives text. What the
customer cannot observe is the quantity of computation the provider spent
producing that text. The provider controls the reasoning budget, the sampling
configuration, the output cap, and any injected system context; none of these is
visible in the response, and all of them change cost.

This is not a hypothetical concern. Reasoning tokens are documented as silently
dropped by some models routed through gateways, and at least one major model
family defaults its reasoning-effort parameter to \texttt{none}, so client code
that reasoned under a previous version silently stops. Measurement studies of
LLM gateways report silent downgrades, model switches, and endpoints that fail
fingerprint verification.

The verification stack that has grown up around hosted inference does not
address this. Each layer binds a different object:

\begin{itemize}
\item \textbf{Environment attestation} (confidential-computing GPUs, TEEs,
transparency-logged enclaves) binds hardware, firmware, and the loaded binary
image.
\item \textbf{Proof of inference} (zkML) binds an output to a declared
architecture under committed weights.
\item \textbf{Transport provenance} (signed API attestations, TLS-transcript
proofs) binds a signature chain over request and response.
\item \textbf{Behavioural auditing} binds the served model's output
distribution to a reference.
\end{itemize}

\noindent Under a \emph{silent reasoning-effort downgrade} --- genuine model,
genuine weights, genuine binary, less computation --- the first three layers
\emph{pass}. The image is unchanged, the weights are genuinely the committed
ones, and the provider signs an honest response to a request it really served.

The Hollow-LLM attack shows the cryptographic layers cannot be repaired to cover
this. A prover may declare a large architecture, commit to weights whose
algebraic structure collapses the effective computation, execute only the small
model, and still satisfy the verifier. The authors name this the \emph{effort
gap} and observe that the natural defences either destroy the privacy the proof
was for, or fall back on hardware trust.

The behavioural literature measures capability directly and so is not subject to
that impossibility. But it answers a different question --- \emph{which model is
this?} --- using stylistic and distributional fingerprints, and its evidence does
not leave the auditor. IRIS, the most advanced work in this line, states the
limitation of its own method: no mechanism allows independent verification, the
audit is auditor-centric, and no cryptographic commitment or public randomness is
used.

\paragraph{Contribution.} We take the gap between these literatures seriously
and close it. Concretely:

\begin{enumerate}
\item \textbf{Semantic effort binding.} We bind effort in the output rather than
in the computation, via task success at the \emph{capability cliff} --- the
difficulty band where the claimed model reliably succeeds and a low-effort
substitute does not. Correctness on a problem requiring $N$ sequential steps is
evidence that $N$ steps were taken. This sidesteps the Hollow-LLM impossibility
by leaving the frame in which it holds.

\item \textbf{Transferable evidence.} We make a behavioural audit publicly
verifiable by committing the challenge pool before a public randomness beacon
selects the sample, and by accumulating evidence with an anytime-valid test
supermartingale. The result is a receipt recomputable offline by anyone.

\item \textbf{Zero provider cooperation.} The protocol requires nothing of the
provider --- no new endpoint, no signature, no attestation, not even awareness
that an audit is occurring.

\item \textbf{Evaluation on production frontier APIs}, with all receipts,
transcripts, and negative results published.
\end{enumerate}

\paragraph{What is not new.} No component of Heartwood is novel in isolation.
Betting martingales and e-values are standard; public randomness beacons,
commit-reveal, and black-box LLM auditing all pre-date this work. The
contribution is the composition and the target: pointing a capability test at
\emph{effort}, and making its output transferable.

\section{Related Work}

\paragraph{The effort gap.} Hollow-LLM formalises the central obstacle: a
zero-knowledge proof certifies consistency with a committed witness under a
declared computation, but that guarantee ``does not inherently bind the amount of
computation actually expended.'' Proposed mitigations---sparsity proofs,
behavioural challenge tests, or retreating to TEEs---are each noted by the
authors as insufficient or as sacrificing the properties the proof was for.
Practical zkML systems remain orders of magnitude too slow for frontier-scale
inference independently of this issue.

\paragraph{Attestation.} Confidential-computing deployments attest hardware,
firmware, and the loaded image; some commercial offerings additionally attest the
served weights. Transparency-logged designs such as Apple's Private Cloud Compute
add verifiable release logs, but serve the operator's own products rather than
third-party API customers. None of these attests computation \emph{per request},
which is what the downgrade attack manipulates.

\paragraph{Signed provenance.} Non-intrusive attestation extensions bind a
request commitment to an output commitment with issuer signatures, and explicitly
disclaim proving ``exact model weights or hardware configuration.''
TLS-transcript proofs can bind a response to a session, but current constructions
are designated-verifier: convincing a third party still requires trusting the
notary.

\paragraph{Behavioural auditing.} Model-substitution detection, rank-based
uniformity testing, and gateway auditing all measure the served distribution
against a reference. These methods are semantically meaningful but, as IRIS
records, produce auditor-centric evidence: a third party cannot verify the audit
without trusting the auditor's frozen plan. Fingerprinting work identifies model
versions from a small number of interactions, again targeting identity rather
than effort.

\paragraph{Anytime-valid inference.} Testing by betting, e-values, and test
supermartingales provide type-I control at arbitrary stopping times. These tools
have been applied to auditing settings, including API evaluation under budget
constraints. We use them off the shelf; our contribution is what we bet on and
what we bind the bet to.

\paragraph{Verified randomness in audits.} Trustless-audit constructions use
public randomness with zero-knowledge proofs to audit training data and model
internals, assuming access to the model. Heartwood targets the black-box API
setting, where no such access exists.

\section{Threat Model}

\paragraph{Parties.} A \emph{provider} serves a hosted model and is untrusted.
An \emph{auditor} runs the protocol and is \emph{also} untrusted --- making the
auditor untrusted is the point of the construction. A \emph{verifier} is any
third party who later checks the receipt. A public randomness \emph{beacon} is
trusted only for unpredictability.

\paragraph{Attack.} The provider serves the claimed model with reduced
computation per request: a lowered or disabled reasoning budget, a tightened
output cap, an injected system instruction suppressing deliberation, or any
combination. The customer's prompt is transmitted unchanged. We stress that
\emph{effort is tokens generated, wherever they appear} --- a configuration that
merely relocates reasoning from a hidden channel to the visible response is not
an effort skim, a point we establish empirically in
Section~\ref{sec:negative}.

\paragraph{Goal.} Produce evidence, checkable by a party who trusts neither the
provider nor the auditor, that the endpoint's capability is materially below the
claimed model's calibrated capability.

\paragraph{Out of scope.} Heartwood binds the \emph{auditor's protocol}, not the
\emph{provider's speech}. A dishonest auditor who fabricates a fully
self-consistent transcript produces a receipt that verifies. Closing this
requires a transport-binding layer, and the two compose cleanly: signed
provenance establishes ``the provider really said this,'' while Heartwood
establishes ``what was said required the compute.'' We demonstrate the
fabrication attack rather than omitting it.

\section{The Heartwood Protocol}

\subsection{Phase order}

The security of the construction rests on the order of five phases. An
implementation that reorders them is not Heartwood.

\begin{enumerate}
\item \textbf{Calibrate.} Estimate the claimed model's success rate on a pool
\emph{disjoint} from the audit pool, and select the capability band. Set $p_0$ to
a conservative lower confidence bound (we use a $99\%$ Wilson bound). The
direction matters: overstating $p_0$ would make honest endpoints look guilty, so
the bound biases the test \emph{against} finding fault.

\item \textbf{Commit.} Publish $C = H(\text{pool})$, a hash over every question
\emph{and} its ground-truth answer, before any query is sent. Binding the answers
prevents an auditor from later re-interpreting what ``correct'' meant.

\item \textbf{Bind to a beacon.} Obtain a public randomness value $R$ published
\emph{after} $C$. The item order is derived from $H(C \parallel R)$.

\item \textbf{Audit.} Query items in the derived order, grade each
deterministically, and accumulate evidence.

\item \textbf{Emit a receipt} containing the seed, commitment, beacon round,
frozen plan, full transcript, and result.
\end{enumerate}

\noindent Committing before the beacon exists is what removes control of the
sample from both parties simultaneously: the auditor cannot select a favourable
subset because $R$ did not exist at commit time, and the provider cannot predict
which requests are audits.

\noindent This guarantee is conditional on $R$ being a genuine chain output, and
that condition is not self-enforcing. Offline verification recomputes the
selection order from the beacon value \emph{carried in the receipt}, so an
auditor who invents $R$ satisfies the check by construction and regains control
of the sample; grinding candidate values until the induced order is favourable
is a form of peeking, and voids the anytime-valid guarantee of
Section~\ref{sec:evidence}. We quantify this in Section~\ref{sec:beacon-anchor}
and specify the anchoring step that closes it.

\subsection{Challenge construction}

Items are drawn from families with machine-checkable ground truth (sequential
state tracking, multi-step arithmetic with an exactness constraint, positional
lookup, modular date arithmetic, inclusion--exclusion counting), parameterised by
a difficulty tier. Grading is deterministic and re-runnable from the response
text alone; any LLM-judge or nondeterministic grader would break verifiability,
since the verifier could not reproduce the graded bit.

Two properties matter. First, items must require genuine \emph{sequential} work,
so that the computation cannot be skipped. Second, and less obviously, difficulty
must be calibrated to the \emph{claimed model}: an item the claimed model cannot
solve carries no evidence, because it fails whether the endpoint is honest or
skimmed. We return to this in Section~\ref{sec:cliff}, where it is the dominant
empirical effect.

\subsection{Deterministic, language-independent derivations}

Both derivations a verifier must reproduce --- the item order and the pool itself
--- are specified without reference to any language's random number generator.
The random source is HMAC-SHA256 in counter mode; indices are drawn by rejection
sampling to avoid modulo bias; the order is Fisher--Yates. Pool items draw from
the same stream keyed by a per-item seed. Cross-language test vectors are
published for both. This is a correctness requirement rather than an aesthetic
one: an earlier version of this work specified only the shuffle and left pool
generation on CPython's Mersenne Twister, which silently made receipts
verifiable only under CPython.

\section{Statistical Guarantee}
\label{sec:evidence}

\subsection{Null hypothesis}

Let $x_t \in \{0,1\}$ indicate whether the endpoint answered item $t$ correctly.
The null is capability-based and one-sided:
\[
H_0 : \; \mathbb{E}[x_t \mid \mathcal{F}_{t-1}] \;\ge\; p_0 ,
\]
that is, \emph{the endpoint is at least as capable as the claimed model}.

\subsection{The bet}

Fix a declared tolerance $p_1 < p_0$: the success rate at or below which the
customer considers the service materially degraded. Define
\[
e_t \;=\; 1 + \lambda\,(p_0 - x_t), \qquad
W_T \;=\; \prod_{t=1}^{T} e_t , \qquad
\lambda \in \Big[0, \tfrac{1}{1-p_0}\Big).
\]

\begin{proposition}
Under $H_0$, $W_T$ is a non-negative supermartingale with $W_0 = 1$.
\end{proposition}

\begin{proof}
$e_t > 0$ for $\lambda < 1/(1-p_0)$, so $W_T \ge 0$. Since $\lambda$ is
predictable,
$\mathbb{E}[e_t \mid \mathcal{F}_{t-1}] = 1 + \lambda(p_0 - \mathbb{E}[x_t \mid
\mathcal{F}_{t-1}]) \le 1$ whenever $\mathbb{E}[x_t \mid \mathcal{F}_{t-1}] \ge
p_0$. Hence $\mathbb{E}[W_T \mid \mathcal{F}_{T-1}] \le W_{T-1}$.
\end{proof}

\begin{theorem}[Anytime validity]
For any stopping time $\tau$ and any $\alpha \in (0,1)$,
$\;\Pr_{H_0}\big[\,\exists\, T : W_T \ge 1/\alpha \,\big] \le \alpha$.
\end{theorem}

\begin{proof}
Immediate from Ville's inequality applied to the non-negative supermartingale
$W_T$ with $W_0 = 1$.
\end{proof}

\noindent This is the property that makes the stopping rule safe to publish. The
auditor may stop whenever they wish --- including after seeing the data, or by
continuing an inconclusive audit later --- without inflating type-I error.
Optional continuation is therefore free, provided the plan
$(p_0, p_1, \alpha, \lambda)$ and the pool are unchanged.

\subsection{Choosing $\lambda$}

Maximising the expected log-growth of $W$ under the alternative $p_1$ gives a
closed form:
\[
\lambda^\star \;=\; \frac{p_0 - p_1}{p_0\,(1 - p_0)} .
\]
$\lambda^\star$ is fixed by the plan before the audit and recomputed by the
verifier from $(p_0, p_1)$; a receipt whose stated $\lambda$ disagrees is
rejected. This is what prevents re-tuning the bet after seeing outcomes.

\subsection{Interpreting a negative result}

A negative result is \emph{not} a certificate of honest service. It means: no
degradation beyond $p_1$ was detected within the queries spent. Sensitivity is
exactly what the auditor declares --- a point we demonstrate empirically in
Section~\ref{sec:negative} rather than merely assert.

\section{Verification}

A verifier performs eight checks, all offline:

\begin{enumerate}\itemsep2pt
\item confirm the receipt is well formed and within resource bounds;
\item regenerate the pool from the seed and confirm the commitment;
\item re-derive the order from the beacon value \emph{carried in the receipt}
and confirm it matches the transcript prefix (this defeats cherry-picking,
reordering, and any beacon substitution that does not also recompute the order
--- but see Section~\ref{sec:beacon-anchor});
\item confirm each response hash;
\item re-grade every response;
\item recompute $\lambda$ from the declared $(p_0,p_1)$;
\item recompute the wealth path;
\item recompute the first crossing and the verdict.
\end{enumerate}

\noindent The verifier contacts neither the auditor nor the provider and
requires no network access. In our reference implementation the entire
verification path uses only the language's standard library.

Check~1 exists because a receipt is hostile input by construction: it is a file
one party hands another. An early version of our verifier raised uncaught
exceptions on thirteen classes of malformed receipt and spent 57 seconds on one
that merely \emph{declared} a five-million-item pool --- a denial of service
against the party doing the checking.

\subsection{What offline verification cannot establish}
\label{sec:beacon-anchor}

Check~3 re-derives the selection order from the beacon the receipt supplies.
The eight checks therefore establish that a receipt is \emph{internally
consistent}; they do not establish that its beacon was ever drawn. An auditor
free to choose $R$ satisfies check~3 by construction and recovers exactly the
control over the sample that the commit-then-beacon ordering was introduced to
remove. Because such an auditor may grind candidate values, this is peeking, and
the guarantee of Section~\ref{sec:evidence} does not apply.

We measured the cost of the attack. Against a 300-item pool and an endpoint
whose true success rate ($0.75$) lies \emph{above} $p_0$ ($0.70$) --- so that no
deficit exists to find --- an uncontrolled beacon accumulates $10^{-57.4}$ of
evidence and correctly returns no deficit. Searching candidate beacon values
found one firing at query~6 after 1{,}017 tries, approximately 0.4~seconds of
computation on one core. The resulting receipt passes all eight offline checks.
The declared $\alpha$ was $0.01$; against a grinding auditor the attained
false-positive rate approaches~1.

The mitigation is to fetch the named round from the chain and confirm the
randomness matches. This requires network access, which is why we state it as a
separate step rather than folding it into offline verification: a verifier
without connectivity must report the beacon as unchecked rather than treat
internal consistency as sufficient. All twelve receipts we publish anchor to
genuine drand rounds. Carrying drand's BLS signature inside the receipt would
permit offline anchoring against the chain public key; we have not implemented
this, and note it as the natural next revision.

We report this because it is the kind of gap a protocol paper is expected to
find in review rather than after deployment, and because the reproduction is
short: it is committed alongside the implementation.

\subsection{Paired audits}

The absolute test requires $p_0$, and where $p_0$ comes from determines what a
receipt is worth. Calibrating through the endpoint under audit is circular: a
provider that skims during calibration depresses $p_0$ to its own current
performance, after which the test cannot fire.

A paired audit avoids $p_0$ entirely. Two endpoints that claim the same model
are queried on the same committed pool in the same beacon-derived order. Items
on which they agree are discarded; among the remainder, the null ``$B$ is no
worse than $A$'' renders an $A$-win and a $B$-win equally likely, so $p_0 = 1/2$
by construction rather than by measurement, and the machinery of
Section~\ref{sec:evidence} applies unchanged. This is McNemar's test in
anytime-valid form.

In simulation at $\alpha = 0.01$ the attained false-positive rate for identical
endpoints is $0.010$; the test is inert when $B$ is superior, fires on 86.7\% of
runs when $B$ is mildly inferior ($0.80$ versus $0.65$), and on all runs when
$B$ is skimmed. It cannot detect two endpoints skimming equally --- no black-box
test can --- but it addresses the deployment most likely to skim in practice: an
intermediary reselling access to an origin it bills as the claimed model.

\section{Evaluation}

\subsection{Setup}

We evaluate against a local 2B model (\texttt{gemma:2b} via Ollama, CPU-only) and
two production frontier endpoints on the Anthropic Messages API
(\texttt{claude-haiku-4-5} and \texttt{claude-opus-5}). In every configuration
the honest and skimmed arms serve \emph{the same model with the same weights};
only the serving configuration differs, and the customer's prompt is transmitted
unchanged. The skimmed arm disables reasoning, injects a system instruction
forbidding working, and caps output.

\subsection{The capability cliff is model-relative}
\label{sec:cliff}

Table~\ref{tab:cliff} reports accuracy by difficulty tier for both frontier
models, 12 items per cell.

\begin{table}[t]
\centering
\begin{tabular}{lcccc}
\toprule
& \multicolumn{2}{c}{Haiku 4.5} & \multicolumn{2}{c}{Opus 5} \\
\cmidrule(lr){2-3}\cmidrule(lr){4-5}
Tier & honest & skimmed & honest & skimmed \\
\midrule
2 & 1.00 & 0.42 & --- & --- \\
3 & 1.00 & 0.00 & 1.00 & \textbf{1.00} \\
5 & 1.00 & 0.00 & 1.00 & 0.00 \\
\bottomrule
\end{tabular}
\caption{Accuracy by difficulty tier. Opus~5 solves an 8-step chain in
\emph{3 output tokens} with no deliberation, so at tier 3 the skimmed endpoint
also scores $1.00$ and every query is uninformative. The cliff appears only at
tier 5.}
\label{tab:cliff}
\end{table}

This is the dominant practical finding. The capability band is a property of the
\emph{claimed model}, not of the protocol. An auditor who reuses a difficulty
tier calibrated for a weaker model obtains a well-formed receipt containing no
information. Calibration is not setup; it is load-bearing.

\subsection{Detection on production frontier APIs}

Table~\ref{tab:frontier} reports live audits at $\alpha = 0.01$.

\begin{table}[t]
\centering
\begin{tabular}{llcrrr}
\toprule
Model & Endpoint & $p_1$ & Queries & Rate & $\log_{10} W$ \\
\midrule
Haiku 4.5 & honest         & 0.30 & 60 & 1.000 & $-0.45$ \\
Haiku 4.5 & \textbf{skimmed} & 0.30 & \textbf{4}  & 0.000 & $+2.64$ \\
Haiku 4.5 & \textbf{50\% dilution} & 0.30 & \textbf{14} & 0.429 & $+2.58$ \\
\midrule
Opus 5 & honest            & 0.30 & 45 & 0.978 & $-0.42$ \\
Opus 5 & \textbf{skimmed}  & 0.30 & \textbf{4}  & 0.000 & $+2.07$ \\
Opus 5 & 50\% dilution     & 0.30 & 45 & 0.600 & $+0.49$ \\
Opus 5 & \textbf{50\% dilution} & 0.65 & \textbf{33} & 0.515 & $+2.04$ \\
\bottomrule
\end{tabular}
\caption{Live audits. Bold rows fired ($W \ge 1/\alpha = 10^2$). No honest arm
produced a false positive; the Haiku honest arm's evidence decayed monotonically
to $10^{-27}$ over 60 queries.}
\label{tab:frontier}
\end{table}

\subsection{The audit is cheaper on more capable models}

Detection cost \emph{falls} as the claimed model becomes more capable:

\begin{center}
\begin{tabular}{lcc}
\toprule
Claimed model & full skim & 50\% dilution \\
\midrule
gemma 2B    & 31 queries & not detected \\
Haiku 4.5   & \textbf{4} & \textbf{14} \\
Opus 5      & \textbf{4} & 33 ($p_1 = 0.65$) \\
\bottomrule
\end{tabular}
\end{center}

\noindent We did not anticipate this. A more capable model fails \emph{harder}
when its computation is removed --- the gap between its honest and skimmed
accuracy widens --- so each query carries more evidence. Verification methods
typically become more expensive as models scale; this one becomes cheaper,
which is favourable precisely on the expensive endpoints a customer most wants
to verify.

\subsection{Style versus capability}

A naive skim that returns short answers is separable by a response-length
classifier alone (AUC $0.941$) and therefore demonstrates nothing. Against a
\emph{padded} skimmer --- one that commits to an answer before deliberating and
then emits a fluent, normal-length justification --- a length classifier
performs at chance (AUC $0.523$) while accuracy separates $0.750$ against
$0.208$.

Style is cheap to fake; capability is not. An adversary can reproduce surface
statistics at negligible cost, but cannot produce a correct answer to a problem
it did not compute. The cheapest way to pass a capability test is to do the
work, which makes the test adversarially robust by construction rather than by
obscurity.

\subsection{Replication and error control}

Across independent audits on the local model, each with its own pool seed and
beacon-derived order, $5/5$ skims were detected at queries
$\{23, 26, 28, 28, 34\}$ (median $28$) and $0/3$ honest endpoints produced false
positives. A Monte Carlo of the same plan predicted a median of $28$ and a 90th
percentile of $42$; theory and practice agreed without tuning. Simulated
false-positive rate at the null boundary was $0.0096$ against a nominal
$\alpha = 0.01$. One honest replication realised a success rate of $0.489$,
marginally above $p_0 = 0.446$; its evidence reached only $10^{+0.10}$, far from
the $10^{2}$ threshold.

\subsection{Two negative results}
\label{sec:negative}

We report both, and publish the receipts.

\paragraph{Disabling a reasoning parameter is not an effort skim.} Our first
frontier configuration disabled extended thinking and capped output at 300
tokens. It produced \emph{no separation whatsoever}: honest and skimmed arms both
scored $1.00$ across tiers 3--5. With room to generate, the model simply reasoned
in the \emph{visible} response instead, at 110--252 output tokens and unchanged
accuracy. This is a display change, not a compute change. Effort is tokens
generated wherever they appear, and a threat model that equates ``reasoning
disabled'' with ``less compute'' is wrong.

\paragraph{Degradation milder than the declared tolerance is not detected.} On
Opus~5, a 50\% dilution realised a mixed success rate of $0.600$ --- milder than
the declared tolerance $p_1 = 0.30$ --- so the bet drifted negative and never
fired within 45 queries. Re-declaring $p_1 = 0.65$ as a separate pre-declared
plan caught the same endpoint in 33 queries. Sensitivity is exactly what the
auditor declares; a milder tolerance costs queries and buys sensitivity.

\subsection{Forgery resistance}

Against real receipts, the verifier catches every in-scope forgery we
constructed: flipping a graded bit, rewriting a response, rewriting a response
together with its hash and grade, dropping items the endpoint passed, reordering
items, substituting a beacon, inflating $p_0$, re-tuning $\lambda$, restating the
verdict, substituting the pool, and marking all failures as successes. As stated
in the threat model, a fully self-consistent fabricated transcript is \emph{not}
caught; we include that attack in the suite as a documented boundary.

\section{Limitations}

\begin{enumerate}\itemsep3pt
\item \textbf{Provider speech is unbound.} Heartwood proves the auditor followed
the protocol, not that the provider produced the responses. Composition with
signed provenance or a TLS-transcript proof is required; current
TLS-transcript constructions are themselves designated-verifier.

\item \textbf{A negative result is weak.} It bounds degradation at the declared
tolerance over the queries spent; it is not a certificate of honest service.

\item \textbf{Behaviourally identical substitution is undetectable.} We argue
behavioural equivalence is the correct equivalence class for an API customer,
but this is a real limit rather than a strength.

\item \textbf{Calibration is required per claimed model} and consumes queries
against the honest endpoint. A miscalibrated tier yields a valid receipt
containing no information.

\item \textbf{Offline verification does not anchor the beacon}, and an auditor
who chooses it manufactures a false verdict in roughly a thousand attempts
(Section~\ref{sec:beacon-anchor}). Anchoring requires network access; offline
anchoring requires carrying the chain signature in the receipt and is not yet
implemented.

\item \textbf{$p_0$ must be anchored outside the audited endpoint.} An endpoint
that skims uniformly, including during calibration, depresses $p_0$ to its own
performance and cannot be caught by an absolute test. This is
information-theoretic rather than an implementation gap: catching it requires a
reference the provider cannot bias. A paired audit removes the dependence on
$p_0$ but substitutes a dependence on a second endpoint believed to serve the
claimed model honestly.

\item \textbf{Indistinguishability is argued, not proven.} Our challenges are
ordinary word problems, a common category of API traffic, but a provider
profiling one customer's distribution could detect a shift. Generating
capability-cliff items in the customer's own domain is future work.

\item \textbf{The skim was configured by us}, not observed in the wild. The
experiments show the audit detects the attack; they do not show any provider
performs it.

\item \textbf{The statistical machinery is standard.} The contribution is the
composition and the threat model, not the test.
\end{enumerate}

\section{Conclusion}

Cryptography cannot bind computational effort, and behavioural audits cannot
transfer. Heartwood bridges the two: it binds effort semantically, at the
capability cliff where an answer is reachable only by spending the computation,
and it makes the resulting evidence transferable by committing the challenge pool
before a public beacon selects the sample and by accumulating evidence with an
anytime-valid supermartingale. The construction requires no cooperation from the
provider, and its receipts are recomputable offline using only a standard
library.

The property we find most promising is the one we did not design for: the audit
grows cheaper as the claimed model grows more capable. Verification usually
scales the wrong way. Here it scales toward exactly the endpoints whose cost
makes verification worth doing.

\paragraph{Artefacts.} The reference implementation, protocol specification,
cross-language test vectors, every receipt reported here (including the two
negative results), and the full adversary suite are public and MIT-licensed at
\url{https://github.com/zcashsensei/heartwood}.

\begin{thebibliography}{99}\small

\bibitem{hollowllm} Hollow-LLM Attack: Computationally Trivial Weights in
Zero-Knowledge Verification of LLM Inference. arXiv:2607.28884, 2026.

\bibitem{iris} Which Model Is Actually Serving You? IRIS: Budgeted Black-Box
Auditing of Model Substitution and Routing Dilution in LLM Gateways.
arXiv:2607.20860, 2026.

\bibitem{rut} Auditing Black-Box LLM APIs with a Rank-Based Uniformity Test.
arXiv:2506.06975, 2025.

\bibitem{payfor} Are You Getting What You Pay For? Auditing Model Substitution
in LLM APIs. arXiv:2504.04715, 2025.

\bibitem{aex} AEX: Non-Intrusive Multi-Hop Attestation and Provenance for LLM
APIs. arXiv:2603.14283, 2026.

\bibitem{llmmap} LLMmap: Fingerprinting for Large Language Models.
arXiv:2407.15847, 2024.

\bibitem{zkllm} zkLLM: Zero-Knowledge Proofs for Large Language Models. 2024.

\bibitem{nanozk} NANOZK: Layerwise Zero-Knowledge Proofs for Verifiable Large
Language Model Inference. arXiv:2603.18046, 2026.

\bibitem{pcc} Apple Security Engineering. Private Cloud Compute: A New Frontier
for AI Privacy in the Cloud. 2024.

\bibitem{zkaudit} Trustless Audits without Revealing Data or Models.
arXiv:2404.04500, 2024.

\bibitem{ramdas} A.~Ramdas, J.~Ville, and collaborators. Hypothesis testing with
e-values, martingales and betting. Tutorial and survey, 2024.

\bibitem{wsr} I.~Waudby-Smith and A.~Ramdas. Estimating means of bounded random
variables by betting. \emph{JRSS-B}, 2023.

\bibitem{ville} J.~Ville. \emph{\'Etude critique de la notion de collectif}.
Gauthier-Villars, 1939.

\bibitem{drand} drand: A Distributed Randomness Beacon Daemon.
\url{https://drand.love}.

\bibitem{tlsn} TLSNotary. Zero-knowledge $\neq$ trustless: what ``publicly
verifiable'' means for zkTLS. 2026.

\bibitem{sok} SoK: Market Microstructure and Verification for Decentralized
Inference. arXiv:2510.15612, 2025.

\end{thebibliography}

\end{document}
